\documentclass[12pt]{article}

\usepackage[portuguese]{babel}
\usepackage{float}
\usepackage{listings}
\usepackage{indentfirst}
\usepackage{amsfonts, amssymb, amsmath}
\usepackage{graphicx}
\usepackage{enumitem}
\usepackage[colorlinks=true, allcolors=blue]{hyperref}
\usepackage[letterpaper,
            top=1.5cm,
            bottom=1.5cm,
            left=2.75cm,
            right=2.75cm,
            marginparwidth=1.75cm]{geometry}

\title{Resolução Lista 1 -- ORD}
\author{Guilherme Lopes}
\date{\today}

\begin{document}

\maketitle
    % \begin{itemize}[label=R:]
    %     \item 
    % \end{itemize}
Supondo um disco com as características abaixo, responda as questões a seguir
\begin{itemize}
    \item 8 superficies
    \item 4096 trilhas/superficie
    \item 110 setores/trilha
    \item 512 bytes/setor
    \item Velocidade de rotação = 5400RPM
    \item Tempo médio de seek = 12ms
    \item Latência média = 5,6ms
\end{itemize}
\begin{enumerate}[label=\arabic*.]
    \item Quantos cilindros o disco possui?
    \begin{itemize}[label=R:]
        \item A quantia de cilindros em um disco é igual ao numero de 
        trilhas/superfície, ou seja, nesse disco temos 4096 cilindros.
    \end{itemize}
    \item Quantos cilindros serão necessários para armazenar um arquivo com 
    80.000 registros de 128 bytes cada?
    \begin{itemize}[label=R:]
        \item Considerando que a capacidade do cilindro é dada por:
        \[\textbf{capacidade das trilhas $\times$ número de trilhas do cilindro}\]
        E a capacidade da trilha é dada por:
        \[\textbf{número de bytes por setor $\times$ número de setores por 
        trilha}\]
        Visto que temos 110 setores por trilha e 512 bytes por setor, a 
        capacidade das trilhas é igual a \[512 \times 110 = 56320\]
        Para armazenar um arquivo de 80.000 registros de 128 bytes cada, 
        precisaremos armazenar 10240000 bytes, por cilindro iremos conseguir 
        armazenar cerca de \[56320 \times 8 = 450560\]
        Para armazenar 10240000 bytes, precisaremos de 
        \[\frac{10240000}{450560} = 22,72\]
        Ou seja, para armazenar um arquivo de 80000 registros de 128 bytes cada, 
        precisaremos de 23 cilindros
        
    \end{itemize}
    \item Sabendo que o tempo de transferência de uma trilha é igual ao tempo 
    de rotação, qual é o tempo (em milissegundos) de transferência de uma trilha?
    \begin{itemize}[label=R:]
        \item Visto que a velocidade de rotação é de 5400 RPM, o tempo necessário 
        (em milissegundos) para fazer uma rotação será dado pela razão 
        \[\frac{5400}{1}=\frac{60000}{x} \Rightarrow 5400x=60000 \Rightarrow 
        x=\frac{60000}{5400} = 11.1ms\]
        Logo, o tempo de transfencia é de 11.1ms
    \end{itemize}
    \item Qual o tempo médio (em milissegundos) estimado para a leitura de uma 
    trilha aleatória do disco
\end{enumerate}
\end{document}